\documentclass[11pt,letterpaper]{article}

\usepackage[margin=1.2cm]{geometry}
\usepackage[utf8]{inputenc}
\usepackage[T1]{fontenc}
\usepackage[spanish]{babel}
\usepackage{amsmath,amssymb}
\usepackage{adjustbox}
\usepackage{tikz}
\usetikzlibrary{arrows.meta,positioning,shapes.geometric,shapes.misc,calc}
\pgfdeclarelayer{bg}
\pgfsetlayers{bg,main}

\pagestyle{empty}

\definecolor{azul}{RGB}{31,78,121}
\definecolor{azulclaro}{RGB}{221,235,247}
\definecolor{verde}{RGB}{226,239,218}
\definecolor{amarillo}{RGB}{255,242,204}
\definecolor{naranja}{RGB}{252,228,214}
\definecolor{gris}{RGB}{242,242,242}

\tikzset{
    terminal/.style={
        rounded rectangle,
        draw=azul,
        fill=azulclaro,
        very thick,
        text width=4.8cm,
        minimum height=1.0cm,
        align=center,
        font=\bfseries\small
    },
    proc/.style={
        rectangle,
        rounded corners=2pt,
        draw=black!65,
        fill=gris,
        thick,
        text width=4.6cm,
        minimum height=1.0cm,
        align=center,
        font=\small
    },
    param/.style={
        rectangle,
        rounded corners=2pt,
        draw=black!70,
        fill=verde,
        thick,
        text width=4.2cm,
        minimum height=0.95cm,
        align=center,
        font=\small
    },
    dec/.style={
        diamond,
        aspect=2.15,
        draw=black!70,
        fill=amarillo,
        thick,
        text width=4.0cm,
        align=center,
        inner sep=1pt,
        font=\small\bfseries
    },
    form/.style={
        rectangle,
        rounded corners=2pt,
        draw=black!60,
        fill=white,
        thick,
        text width=6.4cm,
        minimum height=0.85cm,
        align=center,
        font=\scriptsize
    },
    arrow/.style={
        -{Latex[length=2.5mm]},
        thick,
        draw=black!75
    },
    dashedarrow/.style={
        -{Latex[length=2.5mm]},
        thick,
        dashed,
        draw=black!65
    }
}

\begin{document}

\begin{center}
{\Large\bfseries Diagrama de procedimientos operativos del proceso}\\[2mm]
{\small Ejercicio 2: Análisis de riesgo de crédito y flujo de cálculo por carteras}
\end{center}

\vspace*{\fill}

\noindent\makebox[\textwidth][c]{%
\begin{adjustbox}{center,max totalsize={\textwidth}{0.90\textheight}}
\begin{tikzpicture}[node distance=0.5cm and 1.35cm]

\node[terminal] (start) {Inicio\\Análisis crediticio};

\node[dec, below=0.62cm of start] (sel)
  {Seleccionar 3 carteras\\``Carteras202312.xlsx''};

\node[proc, right=2.2cm of sel] (dates)
  {Elegir 2 fechas\\por cartera};

\node[dec, below=0.72cm of sel] (par)
  {Obtener parámetros\\por cartera y fecha};

\node[param, below left=0.82cm and 1.65cm of par]  (pd)
  {Probabilidad de\\incumplimiento (\(PD\))};
\node[param, below=0.82cm of par] (ead)
  {Exposición al\\incumplimiento (\(EAD\))};
\node[param, below right=0.82cm and 1.65cm of par] (lgd)
  {Severidad de\\pérdida (\(LGD\))};

\node[dec, below=0.88cm of ead] (eli) {Calcular pérdida\\esperada individual};
\node[form, right=2.0cm of eli] (elf)
  {\(EL_i=\mathrm{Exp}_i\times LGD_i\times PD_i\)};

\node[dec, below=0.55cm of eli] (uli) {Calcular pérdida no\\esperada individual};
\node[form, right=2.0cm of uli] (ulf)
  {\(UL_i=\mathrm{Exp}_i\times LGD_i\times \sqrt{PD_i(1-PD_i)}\)};

\node[dec, below=0.55cm of uli] (elp) {Calcular pérdida\\esperada del portafolio};
\node[form, right=2.0cm of elp] (elpf)
  {\(\displaystyle EL_P=\sum_{i=1}^{N}\mathrm{Exp}_i\cdot LGD_i\cdot PD_i\)};

\node[dec, below=0.62cm of elp] (ulp) {Calcular pérdida no\\esperada del portafolio};
\node[form, right=2.0cm of ulp] (ulpf)
  {\(\displaystyle UL_P=\sqrt{\sum_{i=1}^{N}\sum_{j=1}^{N}w_iw_j\,\rho_{ij}\,UL_i\,UL_j}\)};

\node[dec, below=0.62cm of ulp] (car) {Calcular capital\\en riesgo (\(CaR\))};
\node[form, right=2.0cm of car] (carf)
  {\(CaR=EL_P+\alpha\cdot UL_P\)};

\node[dec, below=0.55cm of car] (rep)
  {Repetir con\\\(\rho=20\%\) y \(\rho=40\%\)};

\node[dec, below=0.55cm of rep] (cmp)
  {Comparar fechas,\\correlaciones y riesgo};

\node[terminal, below=0.55cm of cmp] (end) {Fin\\Análisis completado};

\begin{pgfonlayer}{bg}
\draw[arrow] (start) -- (sel);
\draw[arrow] (sel) -- node[left, font=\scriptsize\bfseries]{sí} (par);
\draw[arrow] (sel.east) -- node[above, font=\scriptsize\bfseries]{no} (dates.west);
\draw[dashedarrow] (dates.south) -- ++(0,-0.22) -| (sel.north east);

\draw[arrow] (par.south west) -- (pd.north);
\draw[arrow] (par) -- (ead);
\draw[arrow] (par.south east) -- (lgd.north);

\draw[arrow] (pd.south) -- ++(0,-0.18) -| (eli.west);
\draw[arrow] (ead.south) -- (eli.north);
\draw[arrow] (lgd.south) -- ++(0,-0.18) -| (eli.east);

\draw[arrow] (eli) -- (uli);
\draw[arrow] (uli) -- (elp);
\draw[arrow] (elp) -- (ulp);
\draw[arrow] (ulp) -- (car);
\draw[arrow] (car) -- (rep);
\draw[arrow] (rep) -- (cmp);
\draw[arrow] (cmp) -- (end);

\draw[dashedarrow] (eli.east) -- (elf.west);
\draw[dashedarrow] (uli.east) -- (ulf.west);
\draw[dashedarrow] (elp.east) -- (elpf.west);
\draw[dashedarrow] (ulp.east) -- (ulpf.west);
\draw[dashedarrow] (car.east) -- (carf.west);
\end{pgfonlayer}

\end{tikzpicture}
\end{adjustbox}
}

\vspace*{\fill}

\end{document}
